% Intended LaTeX compiler: pdflatex
\documentclass[a4paper,12pt]{article}
\usepackage[utf8]{inputenc}
\usepackage[T1]{fontenc}
\usepackage{graphicx}
\usepackage{longtable}
\usepackage{wrapfig}
\usepackage{rotating}
\usepackage[normalem]{ulem}
\usepackage{amsmath}
\usepackage{amssymb}
\usepackage{capt-of}
\usepackage{hyperref}
\usepackage{\string~/.emacs.d/common}
\author{mahmood}
\date{\today}
\title{probability homework 2}
\hypersetup{
 pdfauthor={mahmood},
 pdftitle={probability homework 2},
 pdfkeywords={},
 pdfsubject={homework on the subject of conditional probability},
 pdfcreator={Emacs 28.2 (Org mode 9.5.5)}, 
 pdflang={English}}
\begin{document}

\maketitle
see \href{20230127023141-dependence.org}{conditional probability} and \href{20230126225225-probability_tree.org}{probability tree}\\

\begin{question}
consider 3 drawers, the first contains 2 golden coins, the second contains 2 silver coins, the third one silver and one gold, we pick a drawer and remove a coin without looking, given that the coin removed was golden, what is the probability that the remaining coin is also golden?\\
\begin{answer}
\includesvg{/home/mahmooz/brain/out/EKUSh1H}\\

\[
P(\text{remain gold} \mid \text{take gold}) = \frac{P(\text{take gold and remain gold})}{P(\text{take gold})} = \frac{\frac{1}{3}}{\frac{1}{3} + \frac{1}{3} \cdot 0.5} = \frac{2}{3}
\]\\
\end{answer}
\end{question}

\begin{question}
a court convicts a guilty person with probability 0.85, and acquits an innocent person with probability 0.9, 70\% of the people that arrive at the court are guilty\\
\begin{subquestion}
what is the percentage of people that the court convicts?\\
\begin{answer}
\includesvg{/home/mahmooz/brain/out/SscUUcZ}\\

\[
0.3 \cdot 0.1 + 0.7 \cdot 0.85 = 0.625
\]\\
\end{answer}
\end{subquestion}
\begin{subquestion}
given that the court has acquitted a person, what is the probability that the person is guilty?\\
\begin{answer}
\[
P(\text{guilty} \mid \text{acquit}) = \frac{P(\text{acquit guilty})}{P(\text{acquitted})} = \frac{0.7 \cdot 0.15}{0.7 \cdot 0.15 + 0.3 \cdot 0.9} = 0.28
\]\\
\end{answer}
\end{subquestion}
\end{question}

\begin{question}
we roll two dice cubes, what is the probability that one of them shows the number 1 given that both show 2 different numbers after falling?\\
\begin{answer}
\includesvg{/home/mahmooz/brain/out/IDAVY42}\\

\[
P(\text{one} \mid \text{2 different nums}) = \frac{P(\text{2 different nums and one})}{P(\text{2 different nums})} = \frac{2 \cdot \frac{1}{6} \cdot \frac{5}{6}}{\frac{5}{6}} = \frac{1}{3}
\]\\
\end{answer}
\end{question}

\begin{question}
after throwing 10 dice cubes we received atleast 3 ones, what is the probability that we got 3 atleast 3 times?\\
\begin{answer}
refer to \href{probability.org}{probability of atleast}\\
\begin{align*}
  P(\text{atleast three threes} \mid \text{atleast three ones}) &= \frac{P(\text{atleast three threes and atleast three ones})}{P(\text{atleast three ones})}\\
  P(\text{atleast three threes}) &= 1-P(\text{zero threes})-P(\text{one three})-P(\text{two threes})\\
  P(\text{atleast three threes and atleast three ones}) = 1 &- P(\text{zero threes and zero ones})\\
  &- P(\text{zero threes and one one})\\
  &- P(\text{zero threes and two ones})\\
  &- P(\text{one three and zero ones})\\
  &- P(\text{one three and one one})\\
  &- P(\text{one three and two ones})\\
  &- P(\text{two threes and zero ones})\\
  &- P(\text{two threes and one one})\\
  &- P(\text{two threes and two ones})
\end{align*}
\end{answer}
\end{question}

\begin{question}
given 6 jugs, three of which contain 10 black balls and 20 whites, another contains all black, and the last two each contain 85\% white balls and the rest are black, we randomly pick a jug and remove a ball, what is the probability that its white?\\
\begin{answer}
\[
\frac{1}{6} \cdot \left(3 \cdot \frac{20}{30} + 1 + \frac{85}{100}\right)
\]\\
\end{answer}
\end{question}

\begin{question}
an investor buys 2 shares today, the probability of the first stock going up at the end of the day is 0.4 and the second 0.7, the probability of neither going up is 0.2\\
\begin{subquestion}
are the stocks independant?\\
\begin{answer}
we're given the fact \(P(\overline{A \cup B}) = 0.2\)\\
\begin{align*}
  P(A \cup B) &= 1 - P(\overline{A \cup B})\\
  &= 0.8\\
  &= P(A) - P(B) - P(A \cap B)\\
  &= 0.4 + 0.7 - P(A \cap B)\\
  P(A \cap B) &= 0.3\\
  &\neq 0.2
\end{align*}
the events are dependent\\
\end{answer}
\end{subquestion}
\begin{subquestion}
the second day we see that the first one went up, what is the probability that the second has too?\\
\begin{answer}
\[
P(\text{second up} \mid \text{first up}) = \frac{P(\text{second up and first up})}{P(\text{first up})} = \frac{0.3}{0.4}
\]\\
\end{answer}
\end{subquestion}
\begin{subquestion}
the second day we see that the second one didnt go up, what is the probability that the first didnt too?\\
\[
P(\text{first not up} \mid \text{second not up}) = \frac{P(\text{first not up and second not up})}{P(\text{second not up})} = \frac{0.2}{0.3}
\]\\
\end{subquestion}
\end{question}

\begin{question}
in a certain locality its known that 84\% of the residents have driving licenses, 90\% of the men have licenses and 80\% of women do too, we randomly pick a resident and turns out that they have a license, what is the probability that its a woman?\\
\begin{answer}
\[
P(\text{woman} \mid \text{has license}) = \frac{P(\text{woman and has license})}{P(\text{has license})} = \frac{0.5 \cdot 0.8}{0.84} = \frac{0.4}{0.84}
\]\\
\end{answer}
\end{question}
\end{document}